\subsection{Vector bundles and principal bundles}

7.10.1. Let F be a Banach space. A vector bundle M with base B is said to be pure of type F if all the fibers \(M_{b}\) of M (for \(b \in B\) ) are isomorphic (as Banach spaces) to F.

Let M be a vector bundle with base B pure of type F, and let P be the open submanifold of the vector bundle \(\mathcal{L}(\mathrm{F}_{\mathrm{B}}; \mathrm{M})\) composed of pairs \((b, u)\) where \(b \in \mathbf{B}\) and where \(u\) is an isomorphism of \(\mathrm{F}_{b} = \mathrm{F}\) onto \(\mathbf{M}_{b}\) . The group \(\mathbf{GL}(\mathbf{F})\) of automorphisms of F acts on the right on P by setting \((b, u). g = (b, u \circ g)\) for \((b, u) \in \mathbf{P}\) and \(g \in \mathbf{GL}(\mathbf{F})\) . Denote by \(\pi_{\mathbf{P}}\) the map \((b, u) \mapsto b\) from P into B. The quadruple \(\lambda = (\mathbf{P}, \mathbf{GL}(\mathbf{F}), \mathbf{B}, \pi_{\mathbf{P}})\) (where \(\mathbf{GL}(\mathbf{F})\) is endowed with its canonical structure as a group manifold (5.12.2)) is a principal fibration (6.2.1): it is called the frame fibration of M. The map \(((b, u), h) \mapsto u(h)\) from \(\mathbf{P} \times \mathbf{F}\) into M endows M with a structure of fiber bundle associated with \(\lambda\) with typical fiber F (6.5.1).

When \(\mathbf{F} = \mathbf{K}^n\) one can identify an isomorphism \(u\) of \(\mathbf{F}\) onto \(\mathbf{M}_b\) with the basis of \(\mathbf{M}_b\) formed by the image under \(u\) of the canonical basis of \(\mathbf{K}^n\) . The frame bundle of \(\mathbf{M}\) is then identified with the open submanifold of \(\mathbf{M} \times_{\mathbf{B}} \ldots \times_{\mathbf{B}} \mathbf{M}\) formed by the bases \((e_1, \ldots, e_n)\) of the different fibers \(\mathbf{M}_b\) .

Let U be an open subset of B and let \(t = (\mathrm{U}, \varphi, \mathrm{E})\) be a vector chart of M with domain U, with E = F. The map \(b \mapsto (b, t_{b})\) is then a section of P, denoted \(\tilde{t}\) , and the map \(t \mapsto \tilde{t}\) is a bijection from the set of vector charts of M of the form \((\mathrm{U}, \varphi, \mathrm{F})\) onto the set of sections of P over U.

7.10.2. Conversely, let \(\lambda = (\mathbf{Q}, \mathbf{G}, \mathbf{B}, \pi_{\mathbf{Q}})\) be a principal fibration and let \(\varphi\) be a homomorphism of group manifolds from G into the group \(\mathbf{GL}(\mathbf{F})\) of automorphisms of a Banach space F. Let G act on the left on F by setting \(g \cdot h = \varphi(g)(h) (h \in \mathbf{F}, g \in \mathbf{G})\) . Let M be a fiber bundle associated with \(\lambda\) with typical fiber F, and let \(\rho: \mathbf{Q} \times \mathbf{F} \to \mathbf{M}\) be its frame map (6.5.1). Let \(\pi\) be the map from M to B defined by

\[
\pi (\rho (q, h)) = \pi_ {\mathrm{Q}} (q) \quad (q \in \mathbf{Q}, h \in \mathbf{F}).
\]

Let s be a section of Q over an open set U of B; the map

\[
\tilde {s}: (b, h) \mapsto (s (b), h)
\]

is then a bijection from U \(\times\) F onto \(\pi^{-1}(\mathbf{U})\) . On the pair (M, \(\pi\) ) there exists a vector bundle structure with base B, and only one, for which the triples \(t_{s} = (\mathbf{U}, \tilde{s}^{-1}, \mathbf{F})\) are vector bundle charts (for every section s of Q). The manifold structure underlying this structure is that of the associated fiber bundle to \(\lambda\) .

Let \(q \in \mathbf{Q}\) ; set \(b = \pi_{\mathbf{Q}}(q)\) and let \(u\) be the isomorphism of \(\mathbf{F}\) onto \(\mathbf{M}_b\) defined by \(u(h) = \rho(q, h)\) (for \(h \in \mathbf{F}\) ). The map \(f: q \mapsto (b, u)\) is a B-morphism of principal fibrations, compatible with \(\varphi\) , from \((\mathbf{Q}, \mathbf{G}, \mathbf{B}, \pi_{\mathbf{Q}})\) into the frame bundle \((\mathbf{P}, \mathbf{GL}(\mathbf{F}), \mathbf{B}, \pi_{\mathbf{P}})\) of the vector bundle \(\mathbf{M}\) .

7.10.3. Let us resume the notation of Section 7.6. For every \(i \in I\), let

\[
\lambda_ {i} = (\mathbf {P} _ {i}, \mathbf {G} _ {i}, \mathbf {B}, \pi_ {i})
\]

a principal fibration with base B, and suppose that \(G_{i}\) acts on the left on a Banach space \(V_{i}\) by means of a homomorphism

\[
\varphi_{i}: \mathbf{G}_{i} \rightarrow \mathbf{GL}(\mathrm{V}_{i}).
\]

Let \(\mathbf{M}_i\) be a fiber bundle associated with \(\lambda_i\) with typical fiber \(\mathbf{V}_i\) . Set \(\mathcal{M} = (\mathbf{M}_i)_{i\in I}\) and \(\mathcal{V} = (\mathrm{V}_i)_{i\in \mathbf{I}}\) and let \(\lambda\) be the principal product fibration of the \(\lambda_{i}\) over B (6.2.5). Set \(\lambda = (\mathrm{P},\mathrm{G},\mathrm{B},\pi_{\mathbf{P}})\) , with \(\mathbf{G} = \prod_{i\in \mathbf{I}}\mathbf{G}_i\) .

Now let \(\tau\) be a vector functor. For \(g = (g_i) \in G\) , let \(\varphi(g)\) be the element of \(\operatorname{Hom}(\mathcal{V}, \mathcal{V})\) defined by:

\[
\begin{array}{l l} \varphi (g) _ {i} = \varphi_ {i} (g _ {i}) & \text { if } \quad i \in \mathrm{I} _ {+} \\ \varphi (g) _ {i} = \varphi_ {i} (g _ {i}) ^ {- 1} & \text { if } \quad i \in \mathrm{I} _ {-}. \end{array}
\]

The group G then acts on \(\tau(\mathcal{V})\) by means of the morphism \(g\mapsto\tau(\varphi(g))\) from G to \(\mathbf{GL}(\tau(\mathcal{V}))\) .

Let moreover \(x = (x_i)\) be a point of P and let \(b = \pi_{\mathbf{P}}(x)\) . For each \(i\) , the map \(\theta_{x_i}\) defined in no. 6.5.2 is an isomorphism from \(V_i\) onto \((\mathbf{M}_i)_b\) . Let \(\theta_x\) be the element of \(\operatorname{Hom}(\mathcal{V}, ((\mathbf{M}_i)_b)_{i \in I})\) defined by:

\[
\begin{array}{l l} (\theta_ {x}) _ {i} = \theta_ {x _ {i}} & \text { if } \quad i \in I _ {+} \\ (\theta_ {x}) _ {i} = \theta_ {x _ {i}} ^ {- 1} & \text { if } \quad i \in I _ {-}. \end{array}
\]

Let \(\rho\) be the map \((x, h) \mapsto (b, \tau(\theta_x)(h))\) of \(\mathbf{P} \times \tau(\mathcal{V})\) into the vector bundle \(\tau(\mathcal{M})\) ; the map \(\rho\) endows \(\tau(\mathcal{M})\) with a structure of fiber bundle associated with \(\lambda\) with typical fiber \(\tau(\mathcal{V})\) .

These considerations generalize to the case of finite-dimensional vector functors, or of vector functors over a field L equipped with a structure of finite-dimensional K-algebra.

